% Check criterion, deliverable, and the Decision Log question of Day 11.
% The criterion comes from the Day 11 section of the syllabus and from
% Labs/day11/README.md.
\begin{frame}{Check criterion}
  \small
  The instructor checks this at the end of the lab:
  \begin{itemize}
    \item \cmark\ \code{stream\_detect.py} prints \code{Task C1: complete},
          and \code{mjpeg\_rate.py} prints \code{Task D1: complete}.
    \item \cmark\ Live demonstration: the detector draws boxes on the RTSP
          stream of the Raspberry Pi, on the laptop.
    \item \cmark\ The report gives the latency (10 values and the median)
          for two settings of the RTSP stream with the detector, and the bit
          rate and the frame rate of the XIAO stream.
    \item \cmark\ The Decision Log gives numbers and names one trade-off.
  \end{itemize}
  \vskip 0.4em
  \textbf{Hand in:} the file \code{report.md}. Keep the images
  \code{latency\_*.png}: the instructor can ask for one.
\end{frame}

\begin{frame}{Decision Log}
  About 100 words: one design decision, your measured numbers, one
  trade-off.
  \vskip 0.2em
  \begin{exerciseblock}
    A shop wants to count the people at its entrance. One camera hangs at
    the door. A Raspberry Pi 5 stands in the back office, 15 m away, on
    Wi-Fi. The count must change within 1 second after a person enters, and
    no image may leave the shop.
    \begin{itemize}
      \item Where do you run the detector?
      \item Which stream settings do you select?
      \item What leaves the camera, and what leaves the shop?
    \end{itemize}
  \end{exerciseblock}
  \vskip 0.2em
  A good Decision Log:
  \begin{itemize}
    \item gives your measured latency and processed rate for the setting
          that you select, with the method (the clock, 10 values),
    \item compares its bit rate with your Wi-Fi throughput,
    \item names what you lose: for example image quality against
          latency.
  \end{itemize}
\end{frame}
